\documentclass[10.5pt, letterpaper]{article}

% Packages:
\usepackage[
    ignoreheadfoot, % set margins without considering header and footer
    top=0.5 cm, % seperation between body and page edge from the top
    bottom=0.3 cm, % seperation between body and page edge from the bottom
    left=0.8 cm, % seperation between body and page edge from the left
    right=0.8 cm, % seperation between body and page edge from the right
    footskip=0.3 cm, % seperation between body and footer
    % showframe % for debugging 
]{geometry} % for adjusting page geometry
\usepackage{titlesec} % for customizing section titles
\usepackage{tabularx} % for making tables with fixed width columns
\usepackage{array} % tabularx requires this
\usepackage[dvipsnames]{xcolor} % for coloring text
\definecolor{primaryColor}{RGB}{0, 79, 144} % define primary color
\usepackage{enumitem} % for customizing lists
\usepackage{fontawesome5} % for using icons
\usepackage{amsmath} % for math
\usepackage[
    pdftitle={Mustafa Mehmood's CV},
    pdfauthor={Mustafa Mehmood},
    pdfcreator={LaTeX with RenderCV},
    colorlinks=true,
    urlcolor=primaryColor
]{hyperref} % for links, metadata and bookmarks
\usepackage[pscoord]{eso-pic} % for floating text on the page
\usepackage{calc} % for calculating lengths
\usepackage{bookmark} % for bookmarks
\usepackage{lastpage} % for getting the total number of pages
\usepackage{changepage} % for one column entries (adjustwidth environment)
\usepackage{paracol} % for two and three column entries
\usepackage{ifthen} % for conditional statements
\usepackage{needspace} % for avoiding page brake right after the section title
\usepackage{iftex} % check if engine is pdflatex, xetex or luatex

% Ensure that generate pdf is machine readable/ATS parsable:
\ifPDFTeX
    \input{glyphtounicode}
    \pdfgentounicode=1
    % \usepackage[T1]{fontenc} % this breaks sb2nov
    \usepackage[utf8]{inputenc}
    \usepackage{lmodern}
\fi



% Some settings:
\AtBeginEnvironment{adjustwidth}{\partopsep0pt} % remove space before adjustwidth environment
\pagestyle{empty} % no header or footer
\setcounter{secnumdepth}{0} % no section numbering
\setlength{\parindent}{0pt} % no indentation
\setlength{\topskip}{0pt} % no top skip
\setlength{\columnsep}{0cm} % set column seperation




\titleformat{\section}{\needspace{4\baselineskip}\bfseries\large}{}{0pt}{}[\vspace{1pt}\titlerule]

\titlespacing{\section}{
    % left space:
    -1pt
}{
    % top space:
    0.12 cm
}{
    % bottom space:
    0.12 cm
} % section title spacing

\renewcommand\labelitemi{$\circ$} % custom bullet points
\newenvironment{highlights}{
    \begin{itemize}[
        topsep=0.03cm, % Space above the list
        parsep=0.03cm, % Space between paragraphs
        partopsep=0pt, % Additional space when the list starts a new paragraph
        itemsep=0.03cm, % Space between list items
        leftmargin=0.4cm % Left margin
    ]
}{
    \end{itemize}
}

\newenvironment{highlightsforbulletentries}{
    \begin{itemize}[
        topsep=0.10 cm,
        parsep=0.10 cm,
        partopsep=0pt,
        itemsep=0pt,
        leftmargin=10pt
    ]
}{
    \end{itemize}
} % new environment for highlights for bullet entries


\newenvironment{onecolentry}{
    \begin{adjustwidth}{
        0.2 cm + 0.00001 cm
    }{
        0.2 cm + 0.00001 cm
    }
}{
    \end{adjustwidth}
} % new environment for one column entries

\newenvironment{twocolentry}[2][]{
    \onecolentry
    \def\secondColumn{#2}
    \setcolumnwidth{\fill, 4.5 cm}
    \begin{paracol}{2}
}{
    \switchcolumn \raggedleft \secondColumn
    \end{paracol}
    \endonecolentry
} % new environment for two column entries

\newenvironment{header}{
    \setlength{\topsep}{0pt}\par\kern\topsep\centering\linespread{1.15}
}{
    \par\kern\topsep
} % new environment for the header

\newcommand{\placelastupdatedtext}{% \placetextbox{<horizontal pos>}{<vertical pos>}{<stuff>}
  \AddToShipoutPictureFG*{% Add <stuff> to current page foreground
    \put(
        \LenToUnit{\paperwidth-2 cm-0.2 cm+0.05cm},
        \LenToUnit{\paperheight-1.0 cm}
    )
  }%
}%

% save the original href command in a new command:
\let\hrefWithoutArrow\href

% new command for external links:
\renewcommand{\href}[2]{\hrefWithoutArrow{#1}{\ifthenelse{\equal{#2}{}}{ }{#2 }\raisebox{.15ex}{\footnotesize \faExternalLink*}}}


\begin{document}

    
    \newcommand{\AND}{\unskip
        \cleaders\copy\ANDbox\hskip\wd\ANDbox
        \ignorespaces
    }
    \newsavebox\ANDbox
    \sbox\ANDbox{}

    \placelastupdatedtext
    \begin{header}
    \textbf{\fontsize{18 pt}{18 pt}\selectfont Mustafa Mehmood}

    \vspace{0.01 cm}

    \normalsize
    
    \mbox{\hrefWithoutArrow{mailto:mustafamehmood8998@gmail.com}{\color{black}{mustafamehmood8998@gmail.com}}}
    \textbf{\,\textperiodcentered\,}
    {Open to Relocate}
    \textbf{\,\textperiodcentered\,}
    \mbox{\hrefWithoutArrow{tel:+44-7769-504-569}{\color{black}{+44 7769 504 569}}}
    \textbf{\,\textperiodcentered\,}
    \mbox{\hrefWithoutArrow{https://mustafamehmood.vercel.app/}{\color{black}{Website}}}
    \textbf{\,\textperiodcentered\,}
    \mbox{\hrefWithoutArrow{https://www.linkedin.com/in/mustafa-meh}{\color{black}{LinkedIn}}}
    \textbf{\,\textperiodcentered\,}
    \mbox{\hrefWithoutArrow{https://github.com/mustafameh}{\color{black}{GitHub }}}
    
\end{header}

\fontsize{10.2}{11.2}\selectfont

\section{Experience}

\vspace{0.05 cm}

        \begin{twocolentry}{
        \textit{London, UK}   
            
        \textit{June 2025 – Present}}
            \textbf{Applied Research Scientist}
            
            \textit{Thomson Reuters}
        \end{twocolentry}
\vspace{0.05 cm}
        \begin{onecolentry}
            \begin{highlights}
                \item Building agentic systems for CoCounsel, TR's AI legal assistant, that extract risks across M\&A due diligence documents, working closely with lawyers to encode their methodologies, identifying 90\%+ of material risks.
                \item Designed an adversarial evaluation pipeline for agentic solutions with 90\%+ failure detection rate (0.78 Cohen's Kappa vs. expert grades), replacing manual grading and reducing evaluation cycles from days to hours.
                \item Built an orchestrator agent harness and the underlying skill library powering it, comprising modular sub-skills for risk analysis, fact extraction, and document classification, validated through SME evaluations.
                \item Devised context-engineering strategies for documents exceeding context limits: retrieval-then-reason, progressive summarization, and state passed between sub-agents.
                \item Engineered reusable agent tools exposed over MCP for document retrieval, citation verification, cross-reference resolution, and redline comparison, shared across CoCounsel agents.
                \item Fine-tuned LLMs for legal use cases including long-form contract drafting, customization, and summarization.
            \end{highlights}
        \end{onecolentry}

\vspace{0.12 cm}
        
        \begin{twocolentry}{
        \textit{Nottingham, UK}   
            
        \textit{May 2024 – May 2025}}
            \textbf{LLM Data Integrity Specialist}
            
            \textit{Outlier AI}
        \end{twocolentry}
\vspace{0.06 cm}
        \begin{onecolentry}
            \begin{highlights}
                \item Improved tool use in frontier models through RLHF training pipelines, enabling them to chain and execute multi-step actions from high-level user intent, supporting AI development at firms including Google and Meta.
                \item Curated instruction tuning and preference datasets to drive a 20\% improvement in instruction adherence.
            \end{highlights}
        \end{onecolentry}


        
\vspace{0.12 cm}
        \begin{twocolentry}{
        \textit{Selangor, Malaysia}   
            
        \textit{June 2021 - Aug 2023}}
            \textbf{Data Scientist Intern (Part-Time)}
            
            \textit{Department of Neuroscience, University of Nottingham Malaysia}
        \end{twocolentry}
\vspace{0.06 cm}
        
        \begin{onecolentry}
            \begin{highlights}
                \item Engineered feature extraction pipelines to transform raw EEG and eye-tracking clinical data into structured features: coherence, Hjorth parameters, and spectral power bands for EEG; gaze patterns and fixation analysis for eye-tracking.
                \item Built EEG-based diagnostic classifiers for neurological and mental health disorders, reaching 80\% accuracy for Parkinson's and 82\% for Schizophrenia.
                \item For autism screening, represented gaze scanpaths as images and trained a custom transfer-learning computer vision model, benchmarked against a PCA-based classical pipeline, to reach 92\% accuracy.
            \end{highlights}
        \end{onecolentry}  

\section{Education}
\begin{twocolentry}{
    \textit{Nottingham, UK}   
    
    \textit{Sept 2023 – Dec 2024}
    \textit{Sept 2020 – Aug 2023}}
    \textbf{University of Nottingham}

    \textit{\textbf{MSc} Artificial Intelligence - \textbf{Grade}: \textbf{Distinction}}

    \textit{\textbf{BSc} Computer Science \textbf{(Hons)}}
\end{twocolentry}
\vspace{0.04cm}
\begin{onecolentry}
\textbf{\textit{Key Modules}}: Machine Learning, Computer Vision, Natural Language Processing, Big Data, Information Visualisation
\end{onecolentry}

    \section{Key Projects}
    \vspace{0.05 cm}
    \textbf{Sherlock-AI: Character Emulation with LLMs} \hfill {\href{https://agentic-sherlock.vercel.app/}{Website} \href{https://github.com/mustafameh/Sherlock-LLM}{Repo}}

\begin{onecolentry}
    \begin{highlights}
        \item Designed a resource-efficient, transferable character emulation pipeline: fine-tuned a quantized LLaMA model with \textit{LoRA} for Sherlock Holmes emulation, automating dialogue extraction and synthetic data generation from public domain texts.
        
        \item Implemented ReAct-style reasoning for multi-step deductive analysis with dual inference support (OpenRouter API and local GGUF), optimizing for performance and cost across cloud and local setups.

\item Built an interactive story mode where users play as characters in branching Sherlock Holmes mysteries through structured decisions and free-form actions, using templated prompts with streaming structured output to deliver real-time scene generation, configurable writing styles, and auto-generated context summaries.

        
        \item In A/B testing, 72.4\% of users favoured the fine-tuned model over baseline for authenticity, with perfect usability scores.
    \end{highlights}
\end{onecolentry}

  
\vspace{0.2 cm}
\textbf{CourseCompanion: RAG-Powered Course Assistant} \hfill {\href{https://github.com/mustafameh/Course-Content-Q-A}{Repo}}

\begin{onecolentry}
    \begin{highlights}
        \item Developed an educational platform enabling professors to create AI teaching assistants that auto-generate knowledge bases from course materials, featuring chat interfaces, FAQ system, and management dashboards; built with Flask and PostgreSQL using FAISS-backed RAG with LangChain and Flair.
    \end{highlights}
\end{onecolentry}
\vspace{0.2 cm}
\textbf{DeepGaze: Autism Diagnosis with ML and Webcam Eye-Tracking} \hfill {\href{https://youtu.be/c82RrlJVLvo}{Demo Video} \href{https://github.com/mustafameh/Automatic-Autism-Diagnosis-Eyetracking-Machinelearning-Research-Webapplication}{Repo}}

\begin{onecolentry}
    \begin{highlights}
        \item Built a modular Flask web app for webcam-based eye-tracking data collection, visualization, real-time prediction, and model retraining, operationalizing the autism-detection method from my publication.
    \end{highlights}
\end{onecolentry}


\section{Publication}

\noindent
\textbf{Mehmood, M.}, Amin, H. U. (2024). Pre-diagnosis for Autism Spectrum Disorder Using Eye-Tracking and Machine Learning Techniques. \textit{Advances in Brain Inspired Cognitive Systems}, Springer Nature Singapore. \href{https://doi.org/10.1007/978-981-97-1417-9\_23}{}

\vspace{0.15cm}
\section{Technical Skills}   
\begin{onecolentry}
\textbf{Programming Languages:} Python, C, Java, JavaScript, SQL, HTML/CSS
\end{onecolentry}

\vspace{0.05cm}
\begin{onecolentry}
\textbf{Libraries \& Frameworks:}
PyTorch, TensorFlow, Hugging Face (Transformers, trl), scikit-learn, LangChain, LangGraph, Pydantic AI, spaCy, vLLM, FastAPI, Flask, pandas, NumPy, FAISS, Weights \& Biases, Gradio/Streamlit, Claude Agent SDK
\end{onecolentry}

\vspace{0.05cm}
\begin{onecolentry}
\textbf{Tools \& Platforms:}
Git/GitLab CI, Docker, AWS (SageMaker, EC2, Lambda, Bedrock),
Spark/Databricks, PostgreSQL, Pinecone, Runpod, DigitalOcean, MCP, Agent Harness Design, Claude Code, Cursor
\end{onecolentry}


\end{document}
